\chapter{Источник square-root connection: исчерпание существующего меню}

Предыдущий gate показал, что class-D defect требует singular root
connection и pairing condensate. Теперь проверяется, содержится ли нужный
order-four action уже в SU(5), spin, flat-character или Nambu-структурах
проекта.

\section{SU(5)-quarter holonomy}

Используемый в split-sector элемент
\begin{equation}
 h=\exp(i3\pi Y)
\end{equation}
действительно имеет order \(4\) на некоторых Standard Model components.
Но sterile right-handed mode имеет
\begin{equation}
 N^c\sim(\mathbf1,\mathbf1)_0,
 \qquad Y(N^c)=0.
\end{equation}
Она является дополнительным \(SU(5)\)-singlet, отсутствующим в минимальном
\(\mathbf{10}+\overline{\mathbf5}\). Поэтому
\begin{equation}
 h|_{N^c}=1,
 \qquad
 P_5|_{N^c}=h^2|_{N^c}=1.
\end{equation}
SU(5)-quarter branch, разделяющая \(Q,L,T_H\), не создаёт sterile
square-root connection.

\section{Абстрактная gauge-holonomy doublet}

Ранний gauge audit использует U(1)-like charges \(q=\pm1\). При
\(\beta=1/4\) он даёт conjugate phases
\begin{equation}
 (-i,+i),
 \qquad
 (-i)^2=(+i)^2=-1.
\end{equation}
Алгебраически это именно требуемый root pair. Но audit не задаёт
representation map этого U(1)-like charge на neutral sterile sector.
Назначить \(q=\pm1\) постфактум означало бы ввести новый charge assignment.

То же относится к multiplet flat characters. В SU(5) split-table characters
назначены \(\mathbf{10}\), \(\mathbf5\) и \(\mathbf5_H\); sterile singlet в
таблице отсутствует. Character \(\pm i\) можно добавить, но он не выводится
из текущей finite algebra.

\section{Spin и Nambu не являются root-source}

Spin-структуры дают coefficient holonomies \(\pm1\), а не независимую линию
с квадратом \(-1\). Nambu transition
\begin{equation}
 \diag(i,-i)
\end{equation}
переносит basis нулевой моды после того, как pair phase уже изменилась на
\(\pi\). Поэтому он является следствием vortex texture и не может служить
причиной её возникновения. Использовать его как source означало бы
логический круг.

\section{Единственный настоящий root}

Требуемый объект уже классифицирован топологически: это square root
ambient torsion line на
\(\mathbb{RP}^3\setminus\gamma\). Его character на generator равен
\(\pm i\), а meridian holonomy равна \(-1\). Именно поэтому line не
продолжается через core.

Следовательно, root существует математически, но не является smooth field
замороженного action. Он должен войти как один из трёх новых объектов:
\begin{enumerate}
 \item sterile-sector \(Z_4\) flat character;
 \item dynamical \(U(1)\) или \(Z_4\) connection с charge-two pairing field;
 \item disorder operator, включающий nonextendable sector в path integral.
\end{enumerate}

\begin{theorem}[Existing-menu root no-go]
Ни SU(5)-holonomy, ни spin structure, ни Nambu transition не задают
обязательный независимый square-root connection на sterile Majorana mode.
Абстрактная quarter-holonomy doublet содержит правильные phases, но её
sterile charge assignment не выведен. Поэтому defect-route требует явного
расширения finite algebra или topological-sector measure.
\end{theorem}

\begin{protocol}[Definition of Done root-source extension]
Новая root-ветвь допустима только если:
\begin{enumerate}
 \item sterile \(Z_4\)-charge следует из finite algebra и real structure;
 \item root sector включён в path-integral measure до neutrino data;
 \item проверены discrete и global anomaly constraints;
 \item та же representation выводит charge-two pairing field;
 \item pairing Hessian или gap equation создаёт condensate;
 \item extension проходит второй независимый sector test.
\end{enumerate}
\end{protocol}

Текущий итог: нужная фаза в математическом меню существует, но не назначена
sterile-моде никакой обязательной структурой. Поэтому переиспользование
SU(5)-quarter branch не закрывает defect parent action.

Машинный протокол сохранён в
\texttt{s2t\_majorana\_root\_source\_menu\_results.json}.